\section*{Introduction}

Contrary to common intuition, no amount of additional media sampling can recover all four kinetic parameters governing a liver-on-a-chip clearance experiment from media concentration data alone. This is not a limitation of any specific platform, drug, or experimental budget: for the minimal mechanistic structure common to liver-chip clearance analyses, media-only sampling imposes a fundamental information bottleneck that denser sampling, longer incubations, and better fitting algorithms cannot overcome. This paper proves that claim exactly, in closed form, constructs the complete set of parameter combinations it leaves indistinguishable, and then shows that removing this specific bottleneck is still not enough: even once it is removed, the resulting clearance estimate can remain poorly constrained by realistic data, for reasons that have nothing to do with the degeneracy just resolved.

In vitro-to-in vivo extrapolation (IVIVE) of hepatic clearance from conventional assays (microsomes, suspension or plated hepatocytes) systematically underpredicts human pharmacokinetics, by 5- to 10-fold on average across compounds and studies \citep{hallifax2010}. Liver-on-a-chip and related microphysiological systems (MPS) have been proposed as a remedy, offering longer culture viability and more physiologically relevant architecture, and a growing literature applies these platforms to estimate intrinsic clearance (CLint) across several commercial systems, including the CN Bio PhysioMimix \citep{docci2022,tsamandouras2017} and the Javelin liver tissue chip \citep{rajan2023}.

Because liver-chip systems couple medium flow, distribution between medium and cells, metabolism, non-specific binding (NSB) to device materials, and evaporation over multi-day incubations, extracting a clinically interpretable CLint from raw concentration-time data requires more than the single-compartment analysis conventionally applied to microsomal or hepatocyte suspension assays. Several groups have applied formal identifiability analysis to specific configurations. Docci et al.\ \citep{docci2022} performed a priori structural identifiability, via the software DAISY, for three candidate models of their system. Milani et al.\ \citep{milani2022} assessed structural and practical identifiability alongside global sensitivity analysis for a mycophenolate mofetil gut-liver MPS, and Milani et al.\ \citep{milani2024} extended this into a general a priori experimental-design tutorial for organ-on-a-chip systems, demonstrated on midazolam in a gut-liver device, explicitly noting that structural identifiability depends on the specific model structure, dosing, and sampled compartments, and does not guarantee precise practical estimation. Separately, a three-compartment digital-twin approach (DigiLoCS) has been fitted retrospectively to 32 compounds across five published in vitro liver systems, substantially reducing both the bias and the variance of predicted human clearance relative to a conventional depletion analysis \citep{aravindakshan2025}.

Each of these analyses asks whether a specific model, for a specific drug and configuration, is identifiable. None asks the platform-independent question this paper answers: for the minimal mechanistic structure common to liver-chip clearance analyses -- a media compartment exchanging drug with a cell compartment, with first-order NSB and metabolism confined to the cells -- exactly what can the measured media concentration-time curve determine, under the sampling design (media-only, serial) that the overwhelming majority of published studies use \citep{docci2022,rajan2023}, \emph{regardless of which platform or drug generated the data}? That distinction matters: a result proved for one model, one drug, and one device tells a reader whether that experiment was sound; a result proved for the minimal model class itself tells every liver-chip study sharing that structure what it can and cannot claim, before a single data point is collected. And once the answer to that question is known, does resolving it in principle guarantee that the parameters can actually be estimated from realistic, finite, noisy data?

This paper answers both questions. First, using the Laplace transform of the two-compartment system, we prove -- exactly, in closed form, and independently of sampling density, experiment duration, drug identity, or platform -- that media-only sampling recovers only three combinations of the four kinetic parameters (distributional clearance CLd, partition coefficient Kp, intrinsic clearance CLint, and non-specific binding rate kb). We do not stop at showing the fourth combination is unrecoverable: we construct the explicit continuous family of parameter quadruples that are indistinguishable from media data alone, giving any reader a direct, checkable description of the ambiguity rather than an existence statement. We show that independently measuring kb -- for example via a cell-free NSB control, a practice already used on some platforms -- collapses this family to a single point. Second, we test whether resolving this structural degeneracy is sufficient for the parameters to be practically estimable. In doing so we encountered, and report directly, a further finding: standard local methods for quantifying practical identifiability (Fisher-information-based covariance approximations and single-start profile likelihood) are themselves numerically unreliable near this type of structural degeneracy, and only a more computationally demanding multi-start profile likelihood gives a trustworthy answer. That answer shows that fixing kb and evaporation at their true values does not, by itself, sharpen the practical identifiability of CLint under realistic sampling designs. We close by proposing a minimum set of reporting items that would let a reader independently assess both the structural and practical identifiability of a published liver-chip CLint estimate.
